\documentclass[10pt,a4paper]{amsart}
\usepackage[margin=19mm]{geometry}
\usepackage{amsmath,amssymb,mathtools}
\usepackage[scr=rsfs,cal=euler]{mathalfa}
\usepackage{xfrac}
\usepackage[unicode,hidelinks,bookmarks=false]{hyperref}
\newcommand{\ie}{i.e.\ }
\newcommand{\cf}{cf.\ }
\newcommand{\resp}{respectively\ }
\newcommand{\Subsection}{\subsection}
\providecommand{\nobreakdash}{\nobreak\hbox{-}}
\setlength{\emergencystretch}{2em}
\multlinegap0pt
\usepackage{blindtext}
\usepackage{booktabs}
\hypersetup{
    colorlinks=true,
    allcolors=blue
}
\usepackage{nicefrac}
\usepackage{pifont}
\usepackage{tikz}
\usetikzlibrary{arrows.meta}
\usepackage{wrapfig}
\usepackage{cleveref}
\usepackage{mathtools}
\usepackage{booktabs}
\usepackage{nicefrac}
\newtheorem{proposition}{Proposition}
\newtheorem{lemma}{Lemma}
\newtheorem{theorem}{Theorem}
\newtheorem{remark}{Remark}
\newtheorem{example}{Example}
\newcommand{\R}{\mathbb{R}}
\DeclareMathOperator{\Tr}{Tr}
\setlength\parindent{0pt}
\definecolor{src}{RGB}{76,114,176}
\definecolor{tgt}{RGB}{221,132,82}
\newcommand{\cmark}{\ding{51}}
\newcommand{\xmark}{\ding{55}}
\begin{document}
\title[One-Step Generative Surrogate Models via Block-Triangular Jo]{One-Step Generative Surrogate Models via Block-Triangular Joint Drifting}
\author{Nicholas Geissler, Shreya Jha, Ricardo Baptista,; Benjamin Peherstorfer}
\date{}
\maketitle
\noindent\emph{Source and licence.} One-Step Generative Surrogate Models via Block-Triangular Joint Drifting. Version arXiv:2609.26435v1 (CC BY 4.0). This material is distributed under the Creative Commons Attribution 4.0 International licence, http://creativecommons.org/licenses/by/4.0/, as recorded in the arXiv OAI licence metadata for this exact version and shown on the abstract page https://arxiv.org/abs/2609.26435v1. Full rights notice: licenses/arxiv-2609.26435-CC-BY-4.0.txt.\par\smallskip
\noindent\textit{Editorial scope.} Counted unit: the original \texttt{theorem} environment \texttt{thm:ProjectedEquilibriumW2} at source lines 444--459 together with its complete proof at lines 461--695; the delivered document keeps source lines 147--696 so that every referenced label and every citation resolves inside it. Native editable source is the paper's own concatenated TeX; the AMS wrapper is editorial formatting only. No publisher class, upstream script or unrelated artwork is redistributed.
\begin{equation}\label{eq:data}
\mathcal{D} = \{X^{(i)}(t) \,|\, i = 1, \dots, N, t = 0, \dots, T\} \subset \mathbb{R}^n\,, 
\end{equation}
which consists of $i = 1, \dots, N$ trajectories $X^{(i)}(0), \dots, X^{(i)}(T)$ over $t = 0, \dots, T$ time steps.


\subsection{One-step generative surrogate modeling}
We seek a
    map  $g: \R^n \times \R^d \times \mathbb{N} \to \R^n$ such that for $\mu_t$-a.e. $x_t$, we have
    \begin{equation}\label{eq:conditionalmap}
        g(x(t),\cdot,t)_{\sharp}\pi
        =
        p_t(\cdot\,|\,x(t))\,,
    \end{equation}
    where $\pi$ is a suitable reference distribution such as a standard normal. 
The condition \eqref{eq:conditionalmap} implies that if $z\sim\pi$ is a sample from the reference $\pi$, then evaluating the map $g$ at $z$,
    \[
        g(x(t),z,t) \sim p_t(\cdot\,|\,x(t)),
    \]
    gives a sample of the transition law $p_t(\cdot \,|\, x(t))$. 
Correspondingly, we refer to $g$ as a one-step map because a single evaluation $g(x(t),z,t)$ produces a sample from the transition law at the next time point. In particular, no numerical integration of auxiliary dynamics or a sequence of intermediate generative steps are required.



\subsection{Drifting schemes for time marginals}
One approach 
    for learning one-step generators to sample from a target distribution $\eta$ is given by drifting schemes, first introduced in \cite{deng2026drifting}, and further explored in, e.g.,  \cite{han2026wflow,turan2026generativedriftingsecretlyscore,lai2026unifiedviewdriftingscorebased}.
For a parametrized function $f_{\theta}: \R^d  \to \R^n$, where $\theta$ denotes the vector of, e.g., neural-network weights, define the model distribution induced by the pushforward $q_{\theta}=(f_{\theta})_{\sharp}\pi$. 
Drifting schemes iteratively update the parameters $\theta_j$ of $f_{\theta_j}$ over iterations $j = 0, 1, 2, \dots$ such that the model distribution $q_{\theta_j} = (f_{\theta_j})_{\sharp}\pi$ improves the match to the target distribution $\eta$.
The iterative updating is achieved via a drift field $V_{\eta,q_{\theta_j}}: \mathbb{R}^n \to \mathbb{R}^n$, which determines how samples $x \sim q_{\theta_j}$ should be transported, 
    \[
    T_{h,q}(x) = 
    x + h V_{\eta, q}(x)\,,
    \]
    where $h > 0$ is a step size. 
The drifting fields must satisfy $V_{\eta, \eta} = 0$ so that the target distribution $\eta$ is a fixed  point $T_{h,\eta}(x) = x$ of $q \mapsto T_{h,q}$. 
On the parameter level, if at iteration $j$ we have $x = f_{\theta_j}(z)$ with a sample $z \sim \pi$, then the transported sample is 
    \[
    \tilde{x} = T_{h,q_{\theta_j}}(f_{\theta_j}(z)) = f_{\theta_j}(z)  + hV_{\eta, q_{\theta_j}}(f_{\theta_j}(z)).
    \]
Because at iteration $j$ we have $f_{\theta_j}(z) \sim q_{\theta_j}$, the distribution of the transported samples is $\tilde{q}_{j + 1} = (T_{h,q_{\theta_j}})_{\sharp}q_{\theta_j}$.
This defines an iteration: draw samples from the current $q_{\theta_j}$ with $f_{\theta_j}$, transport the samples with the drift field, and then fit $\theta_{j + 1}$ so that $f_{\theta_{j + 1}}$ generates samples close to the transported samples, which is achieved with the loss
    \begin{equation}\label{eq:vanilladriftloss}
        \mathcal{L}(\theta; \theta_j) = \mathbb{E}_{z \sim \pi} \left[\|f_{\theta}(z) - \left(f_{\theta_j}(z)+hV_{\eta,q_{\theta_j}}(f_{\theta_j}(z))\right)\|_2^2 \right].
    \end{equation}
In this manner, each gradient step approximates one fixed-point iteration, with $f_{\theta_j}$ parameterizing the current iterate.


\subsection{Problem formulation: Conditional drifting from trajectory data}
The drifting field $V_{\eta,q_{\theta_j}}$ depends on the target distribution
$\eta$; however, in the data-driven setting, which we consider here, 
the drifting field is constructed from a collection of samples of the target $\eta$ together with samples from the current model
distribution $q_{\theta_j}$. 
In particular, the samples from the target distribution provide the empirical approximation of
$\eta$ used by the drifting procedure.
If we now consider the direct application
of drifting to learn how to sample from the transitional law, this would mean to condition on a state $x(t)$ and take $p_t(\cdot\,|\,x(t))$ to be the target distribution $\eta$.  
The difficulty is that the trajectory data in \eqref{eq:data} do not provide multiple samples from $p_t(\cdot \,|\, x(t))$. Instead, for each observed conditioning state $ x(t) = X^{(i)}(t)$, the training data contain only its single realized successor $X^{(i)}(t+1)$.
Consequently, conditioning the empirical training data on an observed
state $X^{(i)}(t)$ does not provide the drifting procedure with samples that
characterize the corresponding  conditional
$p_t(\cdot\,|\,X^{(i)}(t))$. In particular, just having access to one realization of $X^{(i)}(t + 1)$ for a given $X^{(i)}(t)$
provides no information about the spread, shape, or possible multimodality
of the conditional distribution.
We therefore seek a drifting scheme that can learn the transition law from trajectory data without requiring multiple  samples from the conditional distribution for a given conditioning state.




\section{Block-Triangular Joint Drifting}
We propose to apply drifting to the joint transition law, for which transition-pair samples are available, while restricting the generator to a block-triangular form from which the desired conditional one-step map can be extracted. 


\subsection{Targeting the joint law}
For a fixed time $t$, the training data \eqref{eq:data} contains the transition pairs  $\{(X^{(i)}(t),X^{(i)}(t+1))\}_{i=1}^{N}$, which are samples from the joint law $p_t(\cdot,\cdot)$.
Thus, the training data contains multiple samples from the joint law $p_t(\cdot,\cdot)$ for a fixed $t$ as long as $N > 1$.
Rather than separating the transition pairs at time $t$ via conditioning, we use the joint law $p_t$ as the target distribution of the drifting procedure. In this way, all available
    transition pairs at time $t$ contribute samples to the same drifting target.
We introduce the corresponding parametrized map $
        f_\theta:
        \R^n\times\R^d\times\mathbb{N}
        \to
        \R^{2n}
    $
    and model joint distribution
\begin{equation}\label{eq:TargetJointatTheta}
        q_\theta(t)
        =
        f_\theta(\cdot,\cdot,t)_{\sharp}
        (\mu_t\times\pi).
    \end{equation}
    Note that the model joint distribution depends on time $t$ because $f$ is dependent on time. 

    
\subsection{A block-triangular parametrization}
Applying drifting directly with $f_{\theta}$ while using as target the joint law $p_t$ means that  $f_{\theta}$ produces samples $(X(t),X(t+1))$ jointly, but does not in
    general provide a way to fix an arbitrary current state $x(t)$ and sample
    from the transition law $p_t(\cdot\,|\,x(t))$.
To efficiently sample from the transition law, we consider a specific parametrization of $f_{\theta}$, so that when $f_{\theta}$ has learned to sample from the joint law, we additionally obtain a one-step map for the conditional law.
We therefore follow the block-triangular construction of
\cite{baptista2024conditional} and parametrize $f_{\theta}$ as
    \begin{equation}\label{eq:btidform}
        f_\theta(x(t),z,t)
        =
        \left(
        x(t),
        g_\theta(x(t),z,t)
        \right),
    \end{equation}
    where the conditioning variable $x(t)\sim\mu_t$ is a sample from the time marginal $\mu_t$ and $z\sim\pi$ is a reference sample, which is sampled independently from the conditioning variable.
The parametrization \eqref{eq:btidform}  has a  block triangular form because its first component is independent of $z$. In particular, the first marginal induced by it is $\mu_t$ for every value of $\theta$, while only the second component
    $g_\theta$ depends on $\theta$.
The block-triangular structure ensures that matching the joint law recovers the desired conditional law. In particular,
    if
    \[
        f_\theta(\cdot,\cdot,t)_{\sharp}(\mu_t\times\pi)
        =
        p_t,
    \]
    then
    \[
        g_\theta(x(t),\cdot,t)_{\sharp}\pi
        =
        p_t(\cdot\,|\,x(t))
    \]
    for $\mu_t$-almost every $x(t)$; see Theorem~2.4 of
    \cite{baptista2024conditional}.
Thus, by parametrizing $f_{\theta}$ as in  \eqref{eq:btidform}, it is
    sufficient to train the block-triangular map $f_\theta$ to match the joint
    law $p_t$, and one obtains a one-step map for the transition law via $g_{\theta}$.


\subsection{Drifting field}
\label{subsec:ProjectedWFlow}
We now derive a drifting field that is compatible with the joint law as target and with the block-triangular parametrization as model $f_{\theta}$. 
We begin with the drifting field introduced in \cite{han2026wflow} and applying it to the joint law $p_t$ over the state space $\R^{2n}$, 
\begin{equation}\label{eq:WFlowVelocity}
        V_{p_t,q_{\theta}(t)}^{\varepsilon}(y)
        =
        -
        \nabla_y
        \frac{\delta
        \mathcal{S}_{\varepsilon}
        \big(q_{\theta}(t),p_t\big)}
        {\delta q_{\theta}(t)}(y),
        \qquad
        y\in\R^{2n}\,,
    \end{equation}
where $\mathcal{S}_{\varepsilon}$ denotes the Sinkhorn divergence with entropic regularization parameter $\varepsilon>0$; see, e.g., \cite{pmlr-v89-feydy19a}.  The first variation with respect to the model distribution is denoted as $\delta/\delta q_{\theta}(t)$. 

The field in \eqref{eq:WFlowVelocity} acts on both components of the
    joint state. Our block-triangular parametrization \eqref{eq:btidform},
    however, fixes the first component     and therefore cannot realize motion in the dimensions corresponding to the first block of the state space of the joint law.
 With $y=(y_1,y_2)\in\R^n\times\R^n$, we decompose the joint field \eqref{eq:WFlowVelocity} as
    \[
        V_{p_t,q_{\theta}(t)}^{\varepsilon}(y)
        =
        \left(
            V_1^{\varepsilon}(y),
            V_2^{\varepsilon}(y)
        \right),
    \]
    where $V_1^{\varepsilon},V_2^{\varepsilon}:\R^{2n}\to\R^n$ denote joint field's
    first and second components, respectively.
 We then restrict the joint field to the directions compatible with the block-triangular parametrization. Let $P$ denote the orthogonal projection onto the second block. We define
\begin{equation}\label{eq:ProjectedDrift}
        P V_{p_t,q_{\theta}(t)}^{\varepsilon}(y)
        =
        \left(
            0,
            V_2^{\varepsilon}(y)
        \right),
    \end{equation}
    to obtain the projected field, which is the drifting field that we use to update the conditional component $g_{\theta}$.



\subsection{Loss formulation}
\label{subsec:LossFormulation}

For a fixed time $t$, we insert the projected field
    \eqref{eq:ProjectedDrift} into the drifting loss
    \eqref{eq:vanilladriftloss}, which yields
    \begin{align}
        &\mathcal{L}_t(\theta;\theta_j)
        = 
        \\ &\mathbb{E}_{x(t)\sim\mu_t,z\sim\pi}
        \Big[
        \big\|
            f_{\theta}(x(t),z,t)
            -
            \big(
                f_{\theta_j}(x(t),z,t)
                +
                hP V_{p_t,q_{\theta_j}(t)}^{\varepsilon}
                \big(
                    f_{\theta_j}(x(t),z,t)
                \big)
            \big)
        \big\|_2^2
        \Big]. \nonumber
        \label{eq:ProjectedJointLoss}
    \end{align}
Using the block-triangular parametrization
    \eqref{eq:btidform} and the projected field
    \eqref{eq:ProjectedDrift}, the transported target appearing in
    \eqref{eq:ProjectedJointLoss} can be written as
    \begin{multline}
        f_{\theta_j}(x(t),z,t)
        +
        hP V_{p_t,q_{\theta_j}(t)}^{\varepsilon}
        \big(
            f_{\theta_j}(x(t),z,t)
        \big)
        =\\
        \left(
            x(t),\,
            g_{\theta_j}(x(t),z,t)
            +
            hV_2^{\varepsilon}
            \big(
                f_{\theta_j}(x(t),z,t)
            \big)
        \right).
    \end{multline}
Since $f_{\theta}(x(t),z,t)$ has first component $x(t)$ by
    \eqref{eq:btidform}, the first block of the residual in
    \eqref{eq:ProjectedJointLoss} is identically zero. Therefore,
    \eqref{eq:ProjectedJointLoss} reduces to
    \begin{align}
        \mathcal{L}_t(\theta;\theta_j)
        =
        \mathbb{E}_{x(t)\sim\mu_t,z\sim\pi}
        \left[
        \left\|
            g_{\theta}(x(t),z,t)
            -
            \left(
                g_{\theta_j}(x(t),z,t)
                +
                hV_2^{\varepsilon}
                \big(
                    f_{\theta_j}(x(t),z,t)
                \big)
            \right)
        \right\|_2^2
        \right].
        \label{eq:ConditionalDriftLoss}
    \end{align}
Finally, we average the loss over the available time points,
\begin{equation}\label{eq:ProjectedLoss}
        \mathcal{L}(\theta;\theta_j)
        =
        \mathbb{E}_{t\sim\mathcal{U}(\{0,\ldots,T-1\})}
        \left[
\mathcal{L}_t(\theta;\theta_j)
        \right].
    \end{equation}
For each $t$, the drifting field is constructed from the
    corresponding target $p_t$ and current model distribution
    $q_{\theta_j}(t)$, while the parameters $\theta_j$ are shared across time.




\subsection{Theoretical properties}
\label{sec:CDrift:Theory}
Because the joint field \eqref{eq:WFlowVelocity} is the Sinkhorn
    drifting field considered in \cite{han2026wflow}, we can directly build on
    the well-posedness and other theoretical results established therein.
    In particular, the projected field \eqref{eq:ProjectedDrift} inherits the
    regularity conditions required in 
    \cite{han2026wflow}. Indeed, since $P$ is an orthogonal projection,
    \[
        \|PV_{p_t,q}^{\varepsilon}(y) -PV_{p_t,q}^{\varepsilon}(y')\|_2 \leq \|V_{p_t,q}^{\varepsilon}(y)-V_{p_t,q}^{\varepsilon}(y')\|_2
        \qquad
        \text{for all }y,y'\in\R^{2n},
    \]
    so the growth and Lipschitz bounds satisfied by the joint field are
    preserved under projection. Thus, the well-posedness result of
    \cite{han2026wflow} applies also to the projected field.

We additionally note that $q=p_t$ being a (possibly non-unique) fixed point of $PV_{p_t,q}^{\varepsilon}$ follows from the work in \cite{han2026wflow,pmlr-v89-feydy19a} under suitable assumptions. In particular, it is established in \cite{pmlr-v89-feydy19a} that for distributions with compact support, $q=p_t \implies \mathcal{S}_{\varepsilon}(q,p_t) = 0$, and by regularity and non-negativity of the Sinkhorn divergence, $\mathcal{S}_{\varepsilon}(q,p_t) = 0 \implies V^{\varepsilon}_{p_t,q_{\theta}(t)} \equiv 0$. 
Because $P$ is a projection onto the second component of $V^{\varepsilon}_{p_t,q_{\theta}(t)}$,  $p_t$ remains a fixed-point of the projected velocity field.

For drifting, we would like the reverse direction to hold as well, so that we characterize the drift field as
    \[
       PV_{p_t,q}^{0}=0
    \qquad q\text{-a.e.}
    \qquad\Longleftrightarrow\qquad
    q=p_t
    \]
    within the class of joint distributions whose first marginal is $\mu_t$. This ensures that the projection does not introduce additional zeros of the drift field corresponding to joint distributions different from the target distribution.
This implication is not immediate, because projecting a vector field can eliminate nonzero components. Hence it is possible in principle that the projected field vanishes even though the original field does not. We show that this cannot occur for \eqref{eq:ProjectedDrift} when the model distribution $q$ and the target distribution $p_t$ have the same first marginal $\mu_t$ and the distributions are supported compactly and suitably regular. Importantly, the following statement is restricted to the unregularized drift field \eqref{eq:WFlowVelocity}, i.e., $\varepsilon=0$, and thus the following theorem should be viewed as a statement for an idealized setting. The proof proceeds by representing the unregularized drift field through the quadratic-cost optimal transport map from $q$ to $p_t$. We then show that the condition $PV_{p_t,q}^{0}=0$ forces both components of this transport map to coincide with those of the identity map. Consequently, the optimal transport map is the identity $q$-almost everywhere, which implies $q=p_t$.




\begin{theorem}
\label{thm:ProjectedEquilibriumW2}
Set $\varepsilon = 0$ in the drift field \eqref{eq:WFlowVelocity} and fix a time $t$. Let $p_t,q\in\mathcal P_2(\R^{2n})$ have the same first marginal $\mu_t$, where $\mathcal{P}_2(\R^{2n})$ denotes the space of probability measures over $\R^{2n}$ with finite second moment. 
Assume that $q$ and $p_t$ are supported on the closure of $\Omega_1\times\Omega_2$, where $\Omega_1, \Omega_2 \subset \R^{n}$ are bounded, open, convex sets. Assume further that $q$ and $p_t$ are absolutely continuous with respect to Lebesgue measure and that there exist constants $0<\underline{C}\le \overline{C}<\infty$ such that their densities satisfy Lebesgue-almost everywhere
\begin{equation}\label{eq:BoundedDensities}
    \underline{C}\le q\le \overline{C} \quad\text{ on }\Omega_1\times\Omega_2, \qquad \underline{C}\le p_t\le \overline{C} \quad\text{ on }\Omega_1\times\Omega_2.
\end{equation}
Then
\begin{equation}
\label{eq:ProjectedEquilibriumW2}
    PV_{p_t,q}^{0}=0
    \qquad q\text{-a.e.}
    \qquad\Longleftrightarrow\qquad
    q=p_t.
\end{equation}
\end{theorem}

\begin{proof}
Because $q$ is absolutely continuous with respect to the Lebesgue measure,
Brenier's theorem implies that the quadratic-cost optimal transport from
$q$ to $p_t$ is induced by a map
$T_q^{p_t}:\R^{2n}\to\R^{2n}$ that is unique $q$-almost everywhere and of the form
$T_q^{p_t}=\nabla u$, where $u:\R^{2n}\to\R$ is convex \cite[Theorem~2.26]{Ambrosio2013}. By assumption \eqref{eq:BoundedDensities}, the densities $q$ and $p_t$ are bounded above and bounded away from zero on $\Omega_1 \times \Omega_2$. Therefore, by the regularity theorem for quadratic optimal transport \cite[Theorem~2.27]{Ambrosio2013}, the optimal transport map $T_q^{p_t}$ admits a H\"older-continuous, and hence continuous, representative on $\Omega_1\times\Omega_2$. We use this continuous representative in the following. Furthermore, since $u$ is convex and $\nabla u=T_q^{p_t}$ almost everywhere on $\Omega_1\times\Omega_2$, the continuity of this representative implies that $u$ is continuously differentiable and $\nabla u = T_{q}^{p_t}$ everywhere on $\Omega_1 \times \Omega_2$.




We now relate $T_q^{p_t}$ to the unregularized drifting field used in
\eqref{eq:WFlowVelocity}. For the quadratic cost
$c(y,\bar y)=\frac12\|y-\bar y\|_2^2$, a source Kantorovich potential
associated with the Brenier potential $u$ is
\[
    \Phi_q(y)
    =
    \frac12\|y\|_2^2-u(y),
\]
up to an additive constant; see \cite[Proposition~1.21]{santambrogio2015optimal} and the discussion following that proposition. Now note that 
the lower bound
\eqref{eq:BoundedDensities} implies that the support of $q$ is the closure $\overline{\Omega_1 \times \Omega_2}$, and likewise for $p_t$. Additionally, the quadratic cost belongs
to $C^1(\overline{\Omega_1 \times \Omega_2} \times \overline{\Omega_1 \times \Omega_2})$. Hence the assumptions of
\cite[Proposition~7.18]{santambrogio2015optimal} are satisfied, which shows that the
Kantorovich potential is unique up to an additive constant.



Now recall that for $\varepsilon = 0$, the Sinkhorn divergence $\mathcal{S}_{\varepsilon}$ used in the definition of the drift field \eqref{eq:WFlowVelocity} becomes $\mathcal{S}_0(q, p_t) = \frac{1}{2}W_2^2(q, p_t)$; see, e.g., \cite{han2026wflow}. By \cite[Proposition~7.17]{santambrogio2015optimal}, the Kantorovich potential is a subgradient of the functional $q\mapsto\frac12W_2^2(q,p_t)$ and when it is unique up to additive constants, it represents its first variation, 
\[
\frac{\delta\mathcal S_0(q,p_t)}{\delta q}
=
\Phi_q
\]
up to an additive constant. Consequently, for $q$-almost every
$y\in\Omega_1\times\Omega_2$,
\[
    V_{p_t,q}^0(y)
    =
    -\nabla\Phi_q(y)
    =
    \nabla u(y)-y
    =
    T_q^{p_t}(y)-y.
\]
Hence the unregularized drifting field is precisely the displacement
field of the quadratic optimal transport from $q$ to $p_t$.

Now we are ready to prove the implication
\[
    q=p_t
    \quad\Longrightarrow\quad
    PV_{p_t,q}^0=0
    \quad q\text{-almost everywhere}\,.
\]
Assume that $q=p_t$. Since the identity map transports $q$ to itself
with zero quadratic cost, $W_2(q,q)=0$. Because $T_q^q$ is an optimal
transport map from $q$ to itself,
\[
    \frac12
    \int_{\R^{2n}}
        \|T_q^q(y)-y\|_2^2
        \,q(\mathrm dy)
    =
    \frac12W_2^2(q,q)
    =
    0.
\]
The integrand is nonnegative, and therefore
$T_q^q(y)=y$ for $q$-almost every $y$. It follows from
$V_{q,q}^0=T_q^q-\operatorname{Id}$ that
$V_{q,q}^0=0$ $q$-almost everywhere, and hence
$PV_{q,q}^0=0$ $q$-almost everywhere.

We now prove the converse. Assume that
$PV_{p_t,q}^0=0$ $q$-almost everywhere. The strategy is now to show that under this assumption, the map $T_q^{p_t}$ must be the identity map. We do this component-wise, starting with the second component. Write
$y=(y_1,y_2)\in\R^n\times\R^n$ and decompose the optimal transport map
as
\[
    T_q^{p_t}(y_1,y_2)
    =
    \bigl(
        T_1(y_1,y_2),
        T_2(y_1,y_2)
    \bigr).
\]
Recalling that $P(v_1,v_2)=(0,v_2)$, the identity
$V_{p_t,q}^0=T_q^{p_t}-\operatorname{Id}$ gives
\[
    PV_{p_t,q}^0(y_1,y_2)
    =
    \bigl(
        0,
        T_2(y_1,y_2)-y_2
    \bigr).
\]
Consequently,
\begin{equation}\label{eq:Proof:IdentityT2}
    T_2(y_1,y_2)=y_2
    \qquad
    q\text{-a.e. }(y_1,y_2) \in \Omega_1 \times \Omega_2.
\end{equation}

We next show that \eqref{eq:Proof:IdentityT2} holds everywhere on
$\Omega_1\times\Omega_2$. By \eqref{eq:BoundedDensities},
$q(y)\geq \underline{C} >0$ for Lebesgue-almost every $y \in \Omega_1\times\Omega_2$. Hence, for every measurable
set $A\subset\Omega_1\times\Omega_2$,
\[
    q(A)
    =
    \int_A q(y)\,\mathrm dy
    \geq
    \underline{C}\,|A|,
\]
where $|A|$ denotes the Lebesgue measure of $A$. Consequently,
$q(A)=0$ implies $|A|=0$. Thus, \eqref{eq:Proof:IdentityT2} implies
\[
    T_2(y_1,y_2)=y_2
    \qquad
    \text{for Lebesgue-almost every }(y_1,y_2)
    \in\Omega_1\times\Omega_2.
\]
Since $T_q^{p_t}$ is continuous on $\Omega_1\times\Omega_2$, the map
$(y_1,y_2)\mapsto T_2(y_1,y_2)-y_2$ is continuous there. A continuous
function that vanishes Lebesgue-almost everywhere on an open set must
vanish everywhere on that set. Therefore
\begin{equation}\label{eq:Proof:T2IsIdentity}
    T_2(y_1,y_2)=y_2
    \qquad
    \text{for every }(y_1,y_2)\in\Omega_1\times\Omega_2.
\end{equation}



Let us now consider the first component of $T_q^{p_t}$ and show that it agrees with the identity $q$-almost everywhere. Fix an arbitrary reference point $y_2^0\in\Omega_2$ and define
$h:\Omega_1\to\R$ by
\[
    h(y_1)
    =
    u(y_1,y_2^0)
    -
    \frac12\|y_2^0\|_2^2.
\]
Fix arbitrary $y_1\in\Omega_1$ and $y_2\in\Omega_2$. Since $\Omega_2$
is convex, the line segment
\[
    y_2(r)=(1-r)y_2^0+ry_2,
    \qquad r\in[0,1],
\]
is contained in $\Omega_2$. Define
$\gamma(r)=u(y_1,y_2(r))$. Since $u$ is continuously differentiable, the chain rule gives
\begin{align*}
    \gamma'(r)
    &=
    \nabla_{y_2}u(y_1,y_2(r))
    \cdot
    (y_2-y_2^0) =
    y_2(r)\cdot(y_2-y_2^0),
\end{align*}
where we used
$\nabla_{y_2}u(y_1,y_2)=y_2$ on
$\Omega_1\times\Omega_2$ which follows from \eqref{eq:Proof:T2IsIdentity} with $T_{q}^{p_t} = \nabla u$. Therefore,
\begin{align*}
    u(y_1,y_2)-u(y_1,y_2^0)
    &= \int_0^1 y_2(r)\cdot(y_2-y_2^0)\,\mathrm dr \\
    &= \int_0^1 \bigl((1-r)y_2^0+ry_2\bigr)\cdot(y_2-y_2^0)\,\mathrm dr \\
    &= y_2^0\cdot(y_2-y_2^0)\int_0^1(1-r)\,\mathrm dr
       + y_2\cdot(y_2-y_2^0)\int_0^1r\,\mathrm dr \\
    &= \frac12 y_2^0\cdot(y_2-y_2^0)
       + \frac12 y_2\cdot(y_2-y_2^0) \\
    &= \frac12 (y_2+y_2^0)\cdot(y_2-y_2^0)
     = \frac12\|y_2\|_2^2-\frac12\|y_2^0\|_2^2.
\end{align*}
Hence
\begin{equation}\label{eq:IdentityUhy}
    u(y_1,y_2)
    =
    h(y_1)
    +
    \frac12\|y_2\|_2^2
\end{equation}
for every $(y_1,y_2)\in\Omega_1\times\Omega_2$.
Since $u$ is convex, the restriction $y_1\mapsto u(y_1,y_2^0)$ is convex, and subtracting the constant $\frac12\|y_2^0\|_2^2$ shows that $h$ is convex as well.
Since $u$ is continuously differentiable, $h$ is also continuously differentiable. Differentiating \eqref{eq:IdentityUhy} therefore gives
\[
    T_q^{p_t}(y_1,y_2)
    =
    \nabla u(y_1,y_2)
    =
    \bigl(
        \nabla h(y_1),
        y_2
    \bigr),
\]
everywhere on $\Omega_1 \times \Omega_2$. 

It remains to determine the first component $\nabla h(y_1)$. Let
$P_1:\R^{2n}\to\R^n$ denote the projection
$P_1(y_1,y_2)=y_1$. By assumption, $q$ and $p_t$ have the same first
marginal, so
\[
    (P_1)_\sharp q
    =
    \mu_t
    =
    (P_1)_\sharp p_t.
\]
Since $(T_q^{p_t})_\sharp q=p_t$, we obtain
\begin{align*}
    \mu_t
    &=
    (P_1)_\sharp p_t= (P_1)_\sharp\bigl((T_q^{p_t})_\sharp q\bigr)=
    (P_1\circ T_q^{p_t})_\sharp q\\
    &=
    (\nabla h\circ P_1)_\sharp q=
    (\nabla h)_\sharp\bigl((P_1)_\sharp q\bigr)=
    (\nabla h)_\sharp\mu_t,
\end{align*}
which shows that $\nabla h$ transports $\mu_t$ to itself. We already established that  $h(y_1)=u(y_1,y_2^0)-\frac12\|y_2^0\|_2^2$ is convex. Thus, the map $\nabla h$ is an optimal transport map for the quadratic cost \cite[Theorem~2.13]{Ambrosio2013}. The identity map also transports $\mu_t$ to itself and has zero quadratic cost. Hence the optimal transport cost from $\mu_t$ to itself is zero. Since $\nabla h$ is optimal, it must also attain zero cost. Because the quadratic cost is nonnegative and vanishes only when source and target points coincide, it follows that
\[
    \nabla h(y_1)=y_1
    \qquad
    \mu_t\text{-almost everywhere}.
\]

Since $(P_1)_\sharp q=\mu_t$, the identity $\nabla h(y_1)=y_1$ $\mu_t$-almost everywhere implies, together with the representation $T_q^{p_t}(y_1,y_2)=(\nabla h(y_1),y_2)$, that
\[
    T_q^{p_t}(y_1,y_2) = (y_1,y_2)
\]
for $q$-almost every $(y_1,y_2)\in\Omega_1\times\Omega_2$. Thus $T_q^{p_t}=\operatorname{Id}$ $q$-almost everywhere. Because  $(T_q^{p_t})_\sharp q=p_t$, it follows that
\[
p_t = (T_q^{p_t})_\sharp q = (\operatorname{Id})_\sharp q= q.
\]
This proves the converse implication and hence \eqref{eq:ProjectedEquilibriumW2}.
\end{proof}

\begin{thebibliography}{99}
\bibitem[Ambrosio2013]{Ambrosio2013} L.~Ambrosio and N.~Gigli. \newblock A user's guide to optimal transport. \newblock In {\em Modelling and Optimisation of Flows on Networks: Cetraro, Italy 2009}, volume 2062 of {\em Lecture Notes in Mathematics}, pages 1--155. Springer, 2013.
\bibitem[baptista2024conditional]{baptista2024conditional} R.~Baptista, B.~Hosseini, N.~B. Kovachki, and Y.~M. Marzouk. \newblock Conditional sampling with monotone {GAN}s: From generative models to likelihood-free inference. \newblock {\em SIAM/ASA Journal on Uncertainty Quantification}, 12(3):868--900, 2024.
\bibitem[deng2026drifting]{deng2026drifting} M.~Deng, H.~Li, T.~Li, Y.~Du, and K.~He. \newblock Generative modeling via drifting, 2026.
\bibitem[han2026wflow]{han2026wflow} J.~Han, P.~Li, Q.~Guo, R.~Xu, S.~Ermon, and E.~J. Cand{\`e}s. \newblock One-step generative modeling via wasserstein gradient flows, 2026.
\bibitem[lai2026unifiedviewdriftingscorebased]{lai2026unifiedviewdriftingscorebased} C.-H. Lai, B.~Nguyen, N.~Murata, Y.~Takida, T.~Uesaka, Y.~Mitsufuji, S.~Ermon, and M.~Tao. \newblock A unified view of drifting and score-based models, 2026.
\bibitem[pmlr-v89-feydy19a]{pmlr-v89-feydy19a} J.~Feydy, T.~S{\'e}journ{\'e}, F.-X. Vialard, S.-i. Amari, A.~Trouv{\'e}, and G.~Peyr{\'e}. \newblock Interpolating between optimal transport and {MMD} using sinkhorn divergences. \newblock In {\em Proceedings of the Twenty-Second International Conference on Artificial Intelligence and Statistics}, volume~89 of {\em Proceedings of Machine Learning Research}, pages 2681--2690. PMLR, 2019.
\bibitem[santambrogio2015optimal]{santambrogio2015optimal} F.~Santambrogio. \newblock {\em Optimal Transport for Applied Mathematicians: Calculus of Variations, {PDE}s, and Modeling}, volume~87 of {\em Progress in Nonlinear Differential Equations and Their Applications}. \newblock Birkh{\"a}user, 2015.
\bibitem[turan2026generativedriftingsecretlyscore]{turan2026generativedriftingsecretlyscore} E.~Turan and M.~Ovsjanikov. \newblock Generative drifting is secretly score matching: A spectral and variational perspective, 2026.
\end{thebibliography}
\end{document}
